% !PRINTED 461
\weberpage{461}
\begin{center}
§.\,124.

\gesperrt{Das volle Invariantensystem.}
\end{center}

Wir können nun nachweisen, dass die Formen $f, \Delta, C, K$ ein
volles Invariantensystem der Gruppe sind, d.\,h. dass alle Invarianten
der Gruppe als ganze rationale Functionen von diesen vier dargestellt
werden können. Wir stützen uns dabei auf das Theorem von
\gesperrt{Bezout} (Bd.\,I, §.\,49), dass zwei Curven von der $m^{\text{ten}}$
und $n^{\text{ten}}$ Ordnung, die mehr als $mn$ Punkte gemeinsam haben,
einen gemeinsamen Curventheil haben müssen. Berührungspunkte sind
dabei als doppelt oder mehrfach zu zählen. Alle Schnittpunkte zweier
Invariantencurven bilden entweder ein verbundenes System, oder sie
zerfallen in Systeme verbundener Punkte, und alle Punkte eines
solchen Systemes sind gleich oft zu zählen.

Es sei $\Phi$ eine Invariante $m^{\text{ter}}$ Ordnung, die die Function
$f$ nicht als Factor enthält, und unter den Schnittpunkten der Curve
$\Phi$ und $f$ mögen $h_1$ mal die Pole $P_7$, $h_2$ mal die Pole $P_3$,
$h_3$ mal die Pole $\overline{P}_2$ vorkommen. Ausserdem sollen noch $h_4$
Systeme von je 168 verbundenen Punkten vorkommen, von denen auch
(bei Berührung) mehrere Systeme in ein mehrfach gezähltes zusammenfallen
können. Dann ist die Anzahl aller Schnittpunkte beider Curven
\[
4m = 24h_1 + 56h_2 + 84h_3 + 168h_4,
\]
und daher
\weq{1}{m = 6h_1 + 14h_2 + 21h_3 + 42h_4,}
worin $h_1, h_2, h_3, h_4$ nicht negative ganze Zahlen bedeuten, die
auch nicht alle verschwinden können. Der kleinste Werth, den $m$
haben kann, ist daher 6, und eine Invariantencurve $6^{\text{ter}}$
Ordnung muss durch die 24 Punkte $P_7$ gehen, wie wir es von $\Delta$
schon nachgewiesen haben. Ist $\Delta'$ eine zweite Invariante
$6^{\text{ter}}$ Ordnung, die also auch durch die Punkte $P_7$ gehen
muss, so können wir in $\Delta' - a\Delta$ die Constante $a$ so
bestimmen, dass die Invariantencurve $\Delta' - a\Delta = 0$ durch
irgend einen $25^{\text{sten}}$ Punkt von $f$ geht. Dann muss aber, wenn
$\Delta' - a\Delta$ nicht identisch verschwindet, nach dem
\gesperrt{Bezout}'schen Theorem $\Delta' - a\Delta$ durch $f$ theilbar
sein. Der Quotient wäre eine Invariante zweiter Ordnung, die nicht
existirt; folglich muss $\Delta' - a\Delta$ identisch Null sein.
